\newcommand{\tipoRequerimento}{cobrança indevida sobre enquadramento tarifário incorreto}
\newcommand{\tese}{}

\section{DOS FATOS E DOS DIREITOS}

\begin{enumerate}[resume]
    \item Primeiramente, resta evidente a nulidade do laudo de vistoria técnica emitido pela concessionária, uma vez que esta se limitou a alegar genericamente que a carga instalada "não é exclusiva de água, esgoto e saneamento", omitindo-se completamente do dever de descrever, quantificar ou identificar quais seriam as supostas cargas de natureza diversa presentes na unidade consumidora.
    
    \item Ademais, a alegação de que a unidade consumidora alimenta tomadas, ilustrada pela concessionária apenas com a foto do interior de um quadro elétrico, não constitui argumento idôneo para afastar o enquadramento.
    
    \item Isto porque a presença de tomadas de uso geral (TUGs) no interior de painéis elétricos e salas de comando é requisito técnico indispensável para a operação, calibração, instrumentação e manutenção das próprias instalações de saneamento, servindo estritamente ao suporte da atividade principal da unidade.
    
    \item Portanto, visto que tais circuitos terminais são cargas auxiliares e essenciais ao funcionamento do sistema de saneamento básico, a mera existência de tomadas não desnatura a destinação exclusiva da infraestrutura e nem afasta o direito ao benefício tarifário pleiteado.
    
    \item Contudo, a mera invocação genérica de critérios restritivos e sem lastro técnico descritivo por parte da concessionária viola o dever de fundamentação das decisões e o direito à informação do consumidor, previstos na Resolução Normativa nº 1.000/2021 da ANEEL.
    
    \item Desse modo, visto que a energia elétrica consumida atende estritamente a finalidades vinculadas ao saneamento básico, torna-se indevida e sem amparo regulatório a manutenção da cobrança sob a classificação atual.
    
    \item Por fim, em razão desse enquadramento incorreto e da recusa injustificada da distribuidora em aplicar a tarifa setorial adequada, resta evidente o direito do consumidor à devolução em dobro dos valores pagos a maior, conforme determina a regulação vigente.
\end{enumerate}

\section{DOS PEDIDOS}

\begin{enumerate}[resume]
\item 
Diante do exposto, requer que a Concessionária:

\begin{enumerate}
\item 
Proceda com a correção do enquadramento tarifário da mesma em conformidade com o estabelecido na Resolução Normativa Nº 1.000/2021 da ANEEL.
\item 
Efetue a devolução de todos os valores cobrados a maior, referente aos últimos 10 (dez) anos. Conforme o Despacho Nº 2.006 da ANEEL de 08 de julho de 2024, observando o prazo prescricional de 10 anos, conforme previsto no art. 205 do Código Civil.
\item 
Devolva os valores cobrados a maior em dobro, atualizados monetariamente pelo IPCA e juros de mora de 1\% ao mês calculados pro \textit{rata die}, conformidade com o estabelecido na Resolução Normativa Nº 1.000/2021 da ANEEL.   
\item 
Apresente os argumentos e as provas (comprovações) para não enquadramento (fotos, carga instalada etc).   
\item 
Proceda com a devolução do valor integral, por meio de depósito bancário, na conta corrente do Município, conforme estabelece o § 7º, do art. 323 da REN Nº 1.000/2021 da ANEEL, comunicando por escrito o resultado do deferimento desta reclamação, por meio de e-mail.

\end{enumerate}
\end{enumerate}